\chapter{Research Problem\label{cha:problem}}
% Drafted per POINT-SEQUENCE.md Ch.3 (points 4.1–4.8, re-homed from WP Ch.4).

This chapter states the research problem and the research questions it raises.
It then fixes the scope of the thesis.

\section{Problem Statement\label{sec:goal}}

% [4.1]
NebulaStream keeps all window state in main memory, so a tumbling-window aggregation whose state outgrows its device's memory is killed by the operating system (Section~\ref{sec:bg-metrics}).
Storage bandwidth and access latency are orders of magnitude worse than memory's, so moving state to storage is not a free substitute to provision more memory.
The problem of this thesis is how to store the tumbling-window operator state beyond memory in NebulaStream despite that gap.
The selected state backend is RocksDB, configured in a research-informed manner: the three parameters that prior work identifies as decisive under a memory budget, the compression, the write buffer size, and the block cache size (Section~\ref{sec:bg-rocksdb}), are fixed by measurement before the evaluation.
It must preserve query results and work inside the push-based execution model and its shared global buffer pool (Section~\ref{sec:bg-nes}).

\section{Research Questions\label{sec:rqs}}

The design and evaluation answer six research questions.

% [4.2]
\textbf{RQ0 (design).}
How can an Out-of-Core storage solution be integrated into NebulaStream’s push-based execution model?
The design and implementation in Chapter~\ref{cha:solution} answer this question by construction.

% [4.3]
\textbf{RQ1 (effectiveness).}
To what extent can an Out-of-Core RocksDB based backend enable NebulaStream to process tumbling-window workloads whose state exceeds available memory?

% [4.4]
\textbf{RQ2 (cost).}
How does a research-informed RocksDB configuration perform in terms of throughput and tail latency when used as an Out-of-Core state backend for tumbling windows in NebulaStream?

% [4.5]
\textbf{RQ3 (memory behavior).}
How does memory usage evolve as state grows past the cap, and what explains any failures that occur?

% [4.6]
\textbf{RQ4 (resource use).}
How does CPU utilization and disk I/O activity evolve as operator state grows beyond memory capacity with an Out-of-Core storage backend in NebulaStream?

% [4.7]
\textbf{RQ5 (storage efficiency).}
How efficiently does RocksDB utilize storage for Out-of-Core operator state in NebulaStream, as reflected in write amplification and storage footprint?


\section{Scope\label{sec:scope}}

% [4.8]
The thesis is scoped deliberately.
It covers tumbling windows only.
It runs on a single node.
It changes neither the query compiler nor the compiled query pipelines.
It excludes fault tolerance, durability guarantees, sliding windows, state migration, and automatic parameter tuning.
Cross-system comparisons are out of scope.
The evaluation probes a deliberately fixed subset of the configuration space, varying only the memory budget, the workload state, the pool size, and the allocation type; its answers to the research questions are conditional on that subset.

The backend succeeds if it completes workloads whose state exceeds physical memory, without OOM failures.
Its results must match those of the in-memory engine when sufficient memory is available.
Efficiency is measured against the same in-memory engine under identical processor constraints without memory limits.

The next chapter presents the backend required for this goal.
